% !TEX root = ../aa_SoM_Chapters.tex

%% HARRIS: Chapter 9
%\xdef\Isectiontitle{Statistical Mechanics} 
%%
%\xdef\IaSubsectiontitle{A Simple Thermodynamic System} 
%\xdef\IbSubsectiontitle{Entropy and Temperature}
%\xdef\IcSubsectiontitle{Einstein's solid}
%\xdef\IdSubsectiontitle{Boltzmann \& Maxwell–Boltzmann distributions}
%\xdef\IeSubsectiontitle{Classical Averages}
%
%\xdef\IfSubsectiontitle{Quantum Distributions} 
%\xdef\IgSubsectiontitle{The Quantum Gas} 
%\xdef\IhSubsectiontitle{Photon Gas}
%\xdef\IiSubsectiontitle{The Laser}
%\xdef\IjSubsectiontitle{Specific Heats} 
%\xdef\IkSubsectiontitle{Debye Model} 

% ö = \%c3\%b6
% ä = \%C3\%A4

\xdef\sectiontitle{\IVsectiontitle} %aiheuttama

%%%%%%%%%%%%%%%
\begin{frame}[label=4_5_a]
\frametitle{\texorpdfstring{\culine[\sectiontitleulinecolor]{\sectiontitle}}{\sectiontitle} }

%\vspace{\chapterTOCvksip}%
\begin{itemize}[leftmargin=0.5cm,itemsep=\chapterTOCitemsep]
    \item[4.1] \hyperlink{4_1_a}{\IVaSubsectiontitle}
    \item[4.2] \hyperlink{4_2_a}{\IVbSubsectiontitle}	
    \item[4.3] \hyperlink{4_3_a}{\IVcSubsectiontitle}	
    \item[4.4] \hyperlink{4_4_a}{\IVdSubsectiontitle}
    \item[4.5] \hyperlink{4_5_a}{\CDAlert[blue]{\IVeSubsectiontitle}}
    \item[4.6] \hyperlink{4_6_a}{\IVfSubsectiontitle}
    \item[4.7] \hyperlink{4_7_a}{\IVgSubsectiontitle}
\end{itemize}

%\endofslide 
\end{frame}
%%%%%%%%%%%%%%%

\xdef\sectiontitle{\IVeSubsectiontitle} 


%%%%%%%%%%%%%%%
\begin{frame}
\frametitle{\texorpdfstring{\culine[\sectiontitleulinecolor]{\sectiontitle}}{\sectiontitle} }
%
\framesubtitle{Studying decay processes}

\begin{itemize}
	\item Nearly all the particles observed are unstable \appear{(excluding proton}\appear{, electron} \appear{\& neutrino)}
	\item They tend to decay \appear{and we can learn from their decay modes:}
	\appear{e.g.~energies}\appear{, masses}\appear{, quantum numbers}
	\item Many of the particles are so unstable and short-lived\appear{, that they can be \CDAlert[fuchsia]{studied only indirectly}}
	\item \CDAlert[blue]{Momentum and energy are always conserved}:
	\appear{if apparent energy is not conserved}\appear{, then mass will be changed}
	\item \CDAlert[red]{Charge is always conserved}
	\item \CDAlert[fuchsia]{Are there some fundamental level laws behind this?}
\end{itemize}


\endofslide 
%\endofsection
\end{frame}
%%%%%%%%%%%%%%%

%%%%%%%%%%%%%%%
\begin{frame}
\frametitle{\texorpdfstring{\culine[\sectiontitleulinecolor]{\sectiontitle}}{\sectiontitle} }
%
\framesubtitle{Studying decay processes}

\begin{itemize}
	\item The \CDAlert[blue]{totalitarian principle} \appear{$\oldrightarrow$ the essence of conservation law:}
	\item[] "Everything not forbidden is compulsory" –Gell-Mann (\& R. Heinlein)%Anything that can happen, does happen"
	\item If a decay or a reaction does not occur\appear{, then there must be a reason for that} \appear{$~\Rightarrow~$ \CDAlert[blue]{Conservation law}}
	
	
	\item Conservation laws in nature are fundamentally linked to \CDAlert[fuchsia]{symmetries}
	
	\item If e.g. the potential function of a particle is constant in the direction of motion\appear{, i.e. in the direction of a spatial coordinate}\appear{, then the potential does not cause forces upon the particle in that direction} \appear{and thus, the linear momentum of the particle will be conserved}
	\implies{ constant potential has translational symmetry (\CDAlert[fuchsia]{translational invariance}) }
	
	\item If no angular torques affect on a body, then its potential will have spherical symmetry etc.

\end{itemize}

%
%\xdef\loc{south}
%\begin{tikzpicture}[remember picture,overlay] % anchor=south
%    \node (A) [yshift=-0.25*\figYshift,xshift=\figXshift, anchor=\loc] at (current page.\loc) {\includegraphics[page=1,width=9.5cm]{Figures_booklet/SOM_11_Fundamental_particles_figures/SOM_Fundamental_particles_Figure_1.png}}; %scale=\figscale. % 8cm max hei
%\end{tikzpicture}

\endofslide 
%\endofsection
\end{frame}
%%%%%%%%%%%%%%%

%%%%%%%%%%%%%%%
\begin{frame}
\frametitle{\texorpdfstring{\culine[\sectiontitleulinecolor]{\sectiontitle}}{\sectiontitle} }
%
\framesubtitle{Emmy Noether}

\begin{itemize}
	\item Similarly, the \CDAlert[red]{conservation of energy} rules out any decay of a particle so that the total mass of the decay products would be \Alert[tunired]{greater than the initial mass} 
	
	\item If kinetic energy increases, then there needs to be a decrease of mass
 \implies{ there is a symmetry linking \CDAlert[fuchsia]{time and energy} with each other}

\item \CDAlert[black!30!fuchsia]{Emmy Noether} presented in 1918 a statement on the conjugate variables in physics. \appear{Each physical symmetry leads to a conservation law}

\item The most common conserving quantities \appear{(in classical physics)} \appear{and the corresponding symmetries are:} \vspace{-1mm}
$$ \small
\begin{array}{rcl}
 \text{Energy (mass)} & \qquad \leftrightarrow & \qquad\text{time reversal symmetry}  \\
 \text{Linear momentum} & \qquad \leftrightarrow &\qquad \text{translational symmetry} \\
 \text{Angular momentum} & \qquad \leftrightarrow  & \qquad\text{rotational symmetry} \\
 \text{Electrical charge} &\qquad \leftrightarrow &  \qquad \text{gauge invariance of electrodynamics} \\
\text{\CDAlert[blue]{Minkowski metric}} &\qquad \leftrightarrow &  \qquad \text{\CDAlert[blue]{special theory of relativity}} 
\end{array} 
$$

\end{itemize}

\endofslide 
%\endofsection
\end{frame}
%%%%%%%%%%%%%%%

%%%%%%%%%%%%%%%
\begin{frame}
\frametitle{\texorpdfstring{\culine[\sectiontitleulinecolor]{\sectiontitle}}{\sectiontitle} }
%
\framesubtitle{Emmy Noether}

\begin{itemize}
	 \item Noether also connected \CDAlert[red]{conjugate variables} to \CDAlert[red]{conservation laws}:
	\begin{itemize}
		 \item Energy–time (($\exp (\rmi E t)$\,))
		\item Momentum–displacement (($\exp (\rmi \vec{k} \cdot \vec{r})$\,))
		\item etc.
	\end{itemize}
	
	\item However, already before \CDAlert[fuchsia]{Noether's theorem}\appear{, these conservation laws were more or less known}\appear{, in the scope of classical physics, based mostly on empirical discoveries }
	
	\item In terms of quantum mechanics\appear{, most of the conservation laws are formed on empirical basis}
	
	\item In many cases, the conservation rule has been found to be obeyed only for a particular type of decay 
	
	\item For instance \CDAlert[blue]{strangeness} \Alert[blue]{is conserved only in strong and electromagnetic interactions}, but not is weak interaction, etc.
\end{itemize}

\endofslide 
%\endofsection
\end{frame}
%%%%%%%%%%%%%%%


%%%%%%%%%%%%%%%
\begin{frame}
\frametitle{\texorpdfstring{\culine[\sectiontitleulinecolor]{\sectiontitle}}{\sectiontitle} }
%
\framesubtitle{Emmy Noether}

\begin{itemize}
	\item The \CDAlert[blue]{conservation laws in quantum mechanics} are connected with certain \CDAlert[blue]{quantum numbers}
	\item Thus, we will present \CDAlert[tuniblue]{new quantum numbers}, which will be conserved, but only in certain selected reactions, processes or interactions: \vspace{1.5mm}
	\begin{itemize} \small
	 \item \CDAlert[red]{Baryon number}: +1 for a baryon, $-1$ for antibaryon, 0 otherwise
	\item \CDAlert[blue]{Lepton number}: one for each lepton generation $L_{e}, L_{\mu}, L_{\tau}$ (new data in 1998)
	\item \CDAlert[fuchsia]{Strangeness}: $-1$ for a strange quark, $+1$ for antistrange, 0 otherwise
	\item \CDAlert[green]{Isospin}
	\item \CDAlert[black!20!antiblue]{Hypercharge}
	\end{itemize}

\end{itemize}


\xdef\loc{south east}
\begin{tikzpicture}[remember picture,overlay] % anchor=south
    \node (A) [yshift=-0.15*\figYshift,xshift=\figXshift, anchor=\loc] at (current page.\loc) {\includegraphics[page=1,width=13cm]{Figures_booklet/SOM_11_Fundamental_particles_figures/SOM_Fundamental_Table_4.png}}; %scale=\figscale. % 8cm max hei
\end{tikzpicture}

\endofslide 
%\endofsection
\end{frame}
%%%%%%%%%%%%%%%


%%%%%%%%%%%%%%%
\begin{frame}
\frametitle{\texorpdfstring{\culine[\sectiontitleulinecolor]{\sectiontitle}}{\sectiontitle} }
%
\framesubtitle{On the lepton number}

\seminormal
\begin{itemize}
	\item According to original Standard Model\appear{, \CDAlert[blue]{neutrinos} are assumed to have no mass, but they are polarized} 
	\item Experiments have show neutrinos to be left-handed\appear{, i.e. their spin is antiparallel with their momentum }
	\item Antineutrinos are right handed\appear{, spin parallel to momentum}

\item In the Standard Model\appear{, mass arises from \CDAlert[fuchsia]{Higgs boson}.} \appear{The same interaction also is assumed to change right-handedness to left-handedness, and vice versa}

\item Thus, since neutrinos have always been found to be left-handed\appear{, this indicates they \CDAlert[tuniblue]{should have no mass}.} \appear{This implies that lepton number is conserved in the weak interaction}\appear{, and, more exactly, each of their \Alert[red]{flavours will be conserved}}

\item Except, ... in 1998 data showed that \CDAlert[blue]{neutrino can change into another}

\end{itemize}

\endofslide 
%\endofsection
\end{frame}
%%%%%%%%%%%%%%%


%%%%%%%%%%%%%%%
\begin{frame}
\frametitle{\texorpdfstring{\culine[\sectiontitleulinecolor]{\sectiontitle}}{\sectiontitle} }
%
\framesubtitle{Feynman diagrams}

\begin{itemize}
	 \item Collisions/reactions of fundamental particles are described by illustrative \CDAlert[red]{Feynman diagrams}

	\item Here we take the convention from our textbook by Harris:
	\begin{itemize}
		 \item Horizontal axis is time
		\item Vertical axis is location
	\end{itemize}
	
	\item Absolute spatial/temporal positions do not matter in the diagrams, rather they help in \CDAlert[red]{bookkeeping of possible particle interactions} 
	
\end{itemize}

% 6.5 cm = midway, 10 cm = 3/4 
\vspace{2.5mm}
\begin{minipage}{9cm}
\begin{itemize}
	\item \CDAlert[blue]{Fermions}, (leptons, hadrons, quarks) are described by a \CDAlert[blue]{straight solid line}, \CDAlert[fuchsia]{bosons} (mediating particles) by a \CDAlert[fuchsia]{wavy line}. A joint in the graph means a reaction
\end{itemize}
\end{minipage}
%\vspace{2.5mm}


\xdef\loc{south east}
\begin{tikzpicture}[remember picture,overlay] % anchor=south
    \node (A) [yshift=-0.25*\figYshift,xshift=\figXshift, anchor=\loc] at (current page.\loc) {\includegraphics[page=1,width=4.5cm]{Figures_booklet/SOM_11_Fundamental_particles_figures/SOM_Fundamental_particles_Figure_10.png}}; %scale=\figscale. % 8cm max hei
\end{tikzpicture}

\endofslide 
%\endofsection
\end{frame}
%%%%%%%%%%%%%%%

%%%%%%%%%%%%%%%
\begin{frame}
\frametitle{\texorpdfstring{\culine[\sectiontitleulinecolor]{\sectiontitle}}{\sectiontitle} }
%
\framesubtitle{Feynman diagrams}

\semismall
\begin{itemize}
	\item Basic rules of the interactions: \vspace{2mm}
	\begin{itemize}[itemsep=1.7mm]
	 \item Only particle possessing \CDAlert[fuchsia]{color} can engage in the \CDAlert[fuchsia]{strong force}
	$\oldrightarrow$ Only fermions involved in the \CDAlert[red]{strong vertex are quarks}
	\item To conserve color the gluon carries antiquark. Note, the quark changes color
	\item All particles can interact electromagnetically, except neutrinos
	\item Weak vertices show that all fundamental fermions may interact via W and Z %bosons
	\item In the \CDAlert[blue]{W case}\Alert[blue]{, charge of quark changes and quarks} \CDAlert[blue]{transmutate}
\end{itemize}

\end{itemize}


\xdef\loc{south east}
\begin{tikzpicture}[remember picture,overlay] % anchor=south
    \node (A) [yshift=-0.25*\figYshift,xshift=\figXshift, anchor=\loc] at (current page.\loc) {\includegraphics[page=1,width=14cm]{Figures_booklet/SOM_11_Fundamental_particles_figures/SOM_Fundamental_particles_Figure_8.png}}; %scale=\figscale. % 8cm max hei
\end{tikzpicture}

\endofslide 
%\endofsection
\end{frame}
%%%%%%%%%%%%%%%

%%%%%%%%%%%%%%%
\begin{frame}
\frametitle{\texorpdfstring{\culine[\sectiontitleulinecolor]{\sectiontitle}}{\sectiontitle} }
%
\framesubtitle{Feynman diagrams}

\begin{itemize}
	\item Fermions \appear{(quarks or leptons)} \appear{can e.g. be destructed to a boson} \appear{or be created from a boson}\appear{, this is described by a connection (a crossing point) in the vertices }
	
	\item Note the time reversal symmetries in the two graphs (Figure 9) on quarks and leptons, respectively

\item The whole process (reaction) can be illustrated by combining the vertices into a more complex diagram, as done in Figure 10
\end{itemize}


\xdef\loc{south west}
\begin{tikzpicture}[remember picture,overlay] % anchor=south
    \node (A) [yshift=-0.25*\figYshift,xshift=-\figXshift, anchor=\loc] at (current page.\loc) {\includegraphics[page=1,width=8.7cm]{Figures_booklet/SOM_11_Fundamental_particles_figures/SOM_Fundamental_particles_Figure_9.png}}; %scale=\figscale. % 8cm max hei
\end{tikzpicture}

\appear{\xdef\loc{south east}%
\begin{tikzpicture}[remember picture,overlay] % anchor=south
    \node (A) [yshift=-0.25*\figYshift,xshift=\figXshift, anchor=\loc] at (current page.\loc) {\includegraphics[page=1,width=4.3cm]{Figures_booklet/SOM_11_Fundamental_particles_figures/SOM_Fundamental_particles_Figure_10.png}};%scale=\figscale. % 8cm max hei
\end{tikzpicture}}

\endofslide 
%\endofsection
\end{frame}
%%%%%%%%%%%%%%%


%%%%%%%%%%%%%%%
\begin{frame}
\frametitle{\texorpdfstring{\culine[\sectiontitleulinecolor]{\sectiontitle}}{\sectiontitle} }
%
\framesubtitle{Feynman diagrams}

11. a) Neutron converts to a proton

\xdef\loc{south east}
\begin{tikzpicture}[remember picture,overlay] % anchor=south
    \node (A) [yshift=-0.25*\figYshift,xshift=\figXshift, anchor=\loc] at (current page.\loc) {\includegraphics[page=1,height=8cm]{Figures_booklet/SOM_11_Fundamental_particles_figures/SOM_Fundamental_particles_Figure_11.png}}; %scale=\figscale. % 8cm max hei
\end{tikzpicture}

\endofslide 
%\endofsection
\end{frame}
%%%%%%%%%%%%%%%


%%%%%%%%%%%%%%%
\begin{frame}
\frametitle{\texorpdfstring{\culine[\sectiontitleulinecolor]{\sectiontitle}}{\sectiontitle} }
%
\framesubtitle{Feynman diagrams}

11. b) Pion decays to a muon. This is the well known decay, in the outer regions of Earths atmosphere the cosmic pions are converted into muons, which then get to Earth within a short time of $2.2 \mu$ s due to Lorenz-contaction of the distance.

%\xdef\loc{south}
%\begin{tikzpicture}[remember picture,overlay] % anchor=south
%    \node (A) [yshift=-0.25*\figYshift,xshift=\figXshift, anchor=\loc] at (current page.\loc) {\includegraphics[page=1,width=9.5cm]{Figures_booklet/SOM_11_Fundamental_particles_figures/SOM_Fundamental_particles_Figure_1.png}}; %scale=\figscale. % 8cm max hei
%\end{tikzpicture}

\endofslide 
%\endofsection
\end{frame}
%%%%%%%%%%%%%%%


%%%%%%%%%%%%%%%
\begin{frame}
\frametitle{\texorpdfstring{\culine[\sectiontitleulinecolor]{\sectiontitle}}{\sectiontitle} }
%
\framesubtitle{Feynman diagrams}

11. c) Sigma-baryon decays to a neutron.
Note, strangeness is not conserved!!

%\xdef\loc{south}
%\begin{tikzpicture}[remember picture,overlay] % anchor=south
%    \node (A) [yshift=-0.25*\figYshift,xshift=\figXshift, anchor=\loc] at (current page.\loc) {\includegraphics[page=1,width=9.5cm]{Figures_booklet/SOM_11_Fundamental_particles_figures/SOM_Fundamental_particles_Figure_1.png}}; %scale=\figscale. % 8cm max hei
%\end{tikzpicture}

\endofslide 
%\endofsection
\end{frame}
%%%%%%%%%%%%%%%


%%%%%%%%%%%%%%%
\begin{frame}
\frametitle{\texorpdfstring{\culine[\sectiontitleulinecolor]{\sectiontitle}}{\sectiontitle} }
%
\framesubtitle{Feynman diagrams}

11. d) Muon decays into a positron

%\xdef\loc{south}
%\begin{tikzpicture}[remember picture,overlay] % anchor=south
%    \node (A) [yshift=-0.25*\figYshift,xshift=\figXshift, anchor=\loc] at (current page.\loc) {\includegraphics[page=1,width=9.5cm]{Figures_booklet/SOM_11_Fundamental_particles_figures/SOM_Fundamental_particles_Figure_1.png}}; %scale=\figscale. % 8cm max hei
%\end{tikzpicture}

\endofslide 
%\endofsection
\end{frame}
%%%%%%%%%%%%%%%


%%%%%%%%%%%%%%%
\begin{frame}
\frametitle{\texorpdfstring{\culine[\sectiontitleulinecolor]{\sectiontitle}}{\sectiontitle} }
%
\framesubtitle{Feynman diagrams}

11. e) Sigma-baryon decays to a lambda-baryon

%\xdef\loc{south}
%\begin{tikzpicture}[remember picture,overlay] % anchor=south
%    \node (A) [yshift=-0.25*\figYshift,xshift=\figXshift, anchor=\loc] at (current page.\loc) {\includegraphics[page=1,width=9.5cm]{Figures_booklet/SOM_11_Fundamental_particles_figures/SOM_Fundamental_particles_Figure_1.png}}; %scale=\figscale. % 8cm max hei
%\end{tikzpicture}

\endofslide 
%\endofsection
\end{frame}
%%%%%%%%%%%%%%%


%%%%%%%%%%%%%%%
\begin{frame}
\frametitle{\texorpdfstring{\culine[\sectiontitleulinecolor]{\sectiontitle}}{\sectiontitle} }
%
\framesubtitle{Feynman diagrams}

11. e) Kaon decays into another kaon and a pion

%\xdef\loc{south}
%\begin{tikzpicture}[remember picture,overlay] % anchor=south
%    \node (A) [yshift=-0.25*\figYshift,xshift=\figXshift, anchor=\loc] at (current page.\loc) {\includegraphics[page=1,width=9.5cm]{Figures_booklet/SOM_11_Fundamental_particles_figures/SOM_Fundamental_particles_Figure_1.png}}; %scale=\figscale. % 8cm max hei
%\end{tikzpicture}

\endofslide 
%\endofsection
\end{frame}
%%%%%%%%%%%%%%%


%%%%%%%%%%%%%%%
\begin{frame}
\frametitle{\texorpdfstring{\culine[\sectiontitleulinecolor]{\sectiontitle}}{\sectiontitle} }
%
\framesubtitle{Example 1}

Consider conservation of charge, energy/mass, angular momentum/spin, strangeness and baryon and lepton numbers, in the following three decays. Indicate whether the suggested decays are possible.

$$
\begin{aligned}
\textbf{I:} \quad && \tau^{-} \quad &\rightarrow \quad \mu^{-} ~+~\nu_{\tau} ~+~\overline{\nu_{\mu}}  &&\\
&&  &  &&\\
\textbf{II:}\quad && \Delta^{+} \quad &\rightarrow \quad \quad n~+~K^{+}  &&\\
&&  &  &&\\
\textbf{III:}\quad && \overline{\Lambda^{o}} \quad &\rightarrow \quad \bar{p}~+~\pi^{+}  &&
\end{aligned}
$$

%$\tau^{-} \rightarrow \mu^{-}+v_{\tau}+\overline{v_{\mu}}$
%$\Xi^{o} \rightarrow K^{+}+K^{-}$
%$\Delta^{+} \rightarrow n+K^{+}$
%$\overline{\Lambda^{o}} \rightarrow \bar{p}+\pi^{+}$

%\xdef\loc{south}
%\begin{tikzpicture}[remember picture,overlay] % anchor=south
%    \node (A) [yshift=-0.25*\figYshift,xshift=\figXshift, anchor=\loc] at (current page.\loc) {\includegraphics[page=1,width=9.5cm]{Figures_booklet/SOM_11_Fundamental_particles_figures/SOM_Fundamental_particles_Figure_1.png}}; %scale=\figscale. % 8cm max hei
%\end{tikzpicture}

\endofslide 
%\endofsection
\end{frame}
%%%%%%%%%%%%%%%


%%%%%%%%%%%%%%%
\begin{frame}
\frametitle{\texorpdfstring{\culine[\sectiontitleulinecolor]{\sectiontitle}}{\sectiontitle} }
%
\framesubtitle{Example 1a } %\& b

$$
\begin{aligned}
\textbf{I:} \quad && \tau^{-} \quad &\rightarrow \quad \mu^{-} ~+~\nu_{\tau} ~+~\overline{\nu_{\mu}}  &&\\
&&  &  &&\\
\text {Energy:} \quad && 1780 \mathrm{\;MeV} \quad&\rightarrow \quad 106 \mathrm{\;MeV}  &&\quad\text {OK} \\
\text {Spin:} \quad&& 1 / 2  \quad&\rightarrow \quad  1 / 2 ~+~ 1 / 2 ~+~ 1 / 2 &&\quad\text {{\sim}OK}
\end{aligned}
$$
\implies{Since Lepton numbers are also conserved, the \CDAlert[red]{decay is allowed} }

\xdef\loc{south east}
\begin{tikzpicture}[remember picture,overlay] % anchor=south
    \node (A) [yshift=-0.25*\figYshift,xshift=\figXshift, anchor=\loc] at (current page.\loc) {\includegraphics[page=1,width=4.5cm]{Figures_booklet/SOM_11_Fundamental_particles_figures/SOM_Fundamental_particles_Figure_12a.png}}; %scale=\figscale. % 8cm max hei
\end{tikzpicture}


%$$
%\tau^{-} \rightarrow \mu^{-}+\nu_{\tau}+\overline{v_{\mu}}
%$$

%$$
%\begin{array}{llcll}
%1780 \mathrm{MeV} & \rightarrow & 106 \mathrm{MeV} & \text { OK } \\
%\text { Lepton numbers conserved } & & \text { OK } \\
%\text { Spin } & 1 / 2 & \rightarrow & 1 / 2 & 1 / 2 & 1 / 2 & \text { OK }
%\end{array}
%$$
%The decay is allowed!
%
%$$
%\begin{array}{lllll}
%1315 \mathrm{MeV} & \rightarrow & 2 \times 494 & \mathrm{MeV} & \text { OK } \\
%\text { Baryon number } \mathrm{B}=+1 & \rightarrow & \mathrm{B}=0 \& \mathrm{~B}=0 & & \underline{\text { not }} \text { OK } \\
%\text { Spin } 1 / 2 & \rightarrow & 0 & 0 & & \text { not OK }
%\end{array}
%$$

%$$
%\begin{aligned}
%&& \Xi^{o} \quad &\rightarrow \quad K^{+} ~+~ K^{-}  &&\\
%&&  &  &&\\
%\text {Energy:} \quad && 1315 \mathrm{\;MeV} \quad&\rightarrow \quad 2\times494 \mathrm{\;MeV}  &&\quad\text {OK} \\
%\text {Baryon number:} \quad&& +1  \quad&\rightarrow\quad 0  ~+~ 0 &&\quad\text {not OK} \\
%\text {Spin:} \quad&& 1 / 2  \quad&\rightarrow \quad  1 / 2 ~+~ 1 / 2 ~+~ 1 / 2 &&\quad\text {{\sim}not OK}
%\end{aligned}
%$$
%\implies{?? the \CDAlert[red]{decay is allowed} }



\endofslide 
%\endofsection
\end{frame}
%%%%%%%%%%%%%%%


%%%%%%%%%%%%%%%
\begin{frame}
\frametitle{\texorpdfstring{\culine[\sectiontitleulinecolor]{\sectiontitle}}{\sectiontitle} }
%
\framesubtitle{Example 1c}

$$
\begin{aligned}
\textbf{II:} \quad && \Delta^{+} \quad &\rightarrow \quad \quad n~+~K^{+}  &&\\
&&  &  &&\\
\text {Energy:} \quad && 1232 \mathrm{\;MeV} \quad&\rightarrow \quad 940 ~+~494 \mathrm{\;MeV}  &&\quad\text {not OK} \\
\text {Baryon number:} \quad&& +1  \quad&\rightarrow\quad +1  ~+~ 0 &&\quad\text {OK} \\
\text {Spin:} \quad&& 3 / 2  \quad&\rightarrow \quad  1 / 2 ~+~ 0 &&\quad\text {{\sim}OK}
\end{aligned}
$$
\implies{Final mass is greater than initial mass, the \CDAlert[blue]{decay is not allowed} }

%$$
%\begin{aligned}
%&& \overline{\Lambda^{o}} \quad &\rightarrow \quad \bar{p}~+~\pi^{+}  &&\\
%&&  &  &&\\
%\text {Energy:} \quad && 1116 \mathrm{\;MeV} \quad&\rightarrow \quad 938 ~+~140 \mathrm{\;MeV}  &&\quad\text {OK } \\
%\text {Baryon number:} \quad&& -1  \quad&\rightarrow\quad -1  ~+~ 0 &&\quad\text {OK } \\
%\text {Spin:} \quad&& 1 / 2  \quad&\rightarrow \quad  1 / 2 ~+~ 0 &&\quad\text {OK }
%\end{aligned}
%$$
%\implies{ All properties are conserved, the \CDAlert[red]{decay is allowed} }


%\xdef\loc{south}
%\begin{tikzpicture}[remember picture,overlay] % anchor=south
%    \node (A) [yshift=-0.25*\figYshift,xshift=\figXshift, anchor=\loc] at (current page.\loc) {\includegraphics[page=1,width=9.5cm]{Figures_booklet/SOM_11_Fundamental_particles_figures/SOM_Fundamental_particles_Figure_1.png}}; %scale=\figscale. % 8cm max hei
%\end{tikzpicture}

\endofslide 
%\endofsection
\end{frame}
%%%%%%%%%%%%%%%


%%%%%%%%%%%%%%%
\begin{frame}
\frametitle{\texorpdfstring{\culine[\sectiontitleulinecolor]{\sectiontitle}}{\sectiontitle} }
%
\framesubtitle{Example 1d}

%$$
%\begin{aligned}
%&& \Delta^{+} \quad &\rightarrow \quad \quad n~+~K^{+}  &&\\
%&&  &  &&\\
%\text {Energy:} \quad && 1232 \mathrm{\;MeV} \quad&\rightarrow \quad 940 ~+~494 \mathrm{\;MeV}  &&\quad\text {not OK} \\
%\text {Baryon number:} \quad&& +1  \quad&\rightarrow\quad +1  ~+~ 0 &&\quad\text {OK} \\
%\text {Spin:} \quad&& 3 / 2  \quad&\rightarrow \quad  1 / 2 ~+~ 0 &&\quad\text {{\sim}OK}
%\end{aligned}
%$$
%\implies{Final mass is greater than initial mass, the \CDAlert[blue]{decay is not allowed} }

$$
\begin{aligned}
\textbf{III:} \quad && \overline{\Lambda^{o}} \quad &\rightarrow \quad \bar{p}~+~\pi^{+}  &&\\
&&  &  &&\\
\text {Energy:} \quad && 1116 \mathrm{\;MeV} \quad&\rightarrow \quad 938 ~+~140 \mathrm{\;MeV}  &&\quad\text {OK } \\
\text {Baryon number:} \quad&& -1  \quad&\rightarrow\quad -1  ~+~ 0 &&\quad\text {OK } \\
\text {Spin:} \quad&& 1 / 2  \quad&\rightarrow \quad  1 / 2 ~+~ 0 &&\quad\text {OK }
\end{aligned}
$$
\implies{ All properties are conserved, the \CDAlert[red]{decay is allowed} }


\xdef\loc{south east}
\begin{tikzpicture}[remember picture,overlay] % anchor=south
    \node (A) [yshift=-0.25*\figYshift,xshift=\figXshift, anchor=\loc] at (current page.\loc) {\includegraphics[page=1,width=4.5cm]{Figures_booklet/SOM_11_Fundamental_particles_figures/SOM_Fundamental_particles_Figure_12b.png}}; %scale=\figscale. % 8cm max hei
\end{tikzpicture}

\endofslide 
%\endofsection
\end{frame}
%%%%%%%%%%%%%%%
